\documentclass[10pt,a4paper]{amsart}
\usepackage[margin=19mm]{geometry}
\usepackage{amsmath,amssymb,mathtools}
\usepackage[scr=rsfs,cal=euler]{mathalfa}
\usepackage{xfrac}
\usepackage[unicode,hidelinks,bookmarks=false]{hyperref}
\newcommand{\ie}{i.e.\ }
\newcommand{\cf}{cf.\ }
\newcommand{\resp}{respectively\ }
\newcommand{\Subsection}{\subsection}
\providecommand{\nobreakdash}{\nobreak\hbox{-}}
\setlength{\emergencystretch}{2em}
\multlinegap0pt
\usepackage{amssymb}
\usepackage{bbm}
\usepackage{tikz}
\usepackage{dsfont}
\newtheorem{proposition}{Proposition}
\newcommand{\kerr}{\text{Ker}}
\title{On $\ell_1$ embeddings of finite metric spaces, and sphere-of-influence graphs}
\author{Stanislav Jabuka, Ehsan Mirbagheri}
\date{Source: arXiv:2512.18975v1 (CC BY 4.0)}
\begin{document}
\maketitle

\noindent\emph{Source and licence.} On $\ell_1$ embeddings of finite metric spaces, and sphere-of-influence graphs. Version arXiv:2512.18975v1 (CC BY 4.0), distributed by arXiv under the Creative Commons Attribution 4.0 International licence, http://creativecommons.org/licenses/by/4.0/, as recorded in the arXiv licence metadata for this exact version and shown on the abstract page https://arxiv.org/abs/2512.18975v1. Full licence notice: \texttt{licenses/arxiv-2512.18975-CC-BY-4.0.txt}.\par\smallskip

\noindent\textit{Editorial scope.} Counted unit: the article's proposition kernel (source0.tex lines 1694--1700). The article's complete original proof of it is delivered with it (source0.tex lines 1702--1756, one \texttt{proof} environment of 55 lines). No mathematics is written, repaired or re-derived: the statement and the proof are the article's own lines, in the article's own notation, and no retained line is substituted or re-flowed.  Retained uncounted as the source-local context the proof needs: lines 1635--1645.  Nothing was omitted from any retained span, so the package carries no omission record.  No retained line contains a citation, so the wrapper carries no bibliography and no bibliography entry.  No retained line contains a figure environment, an image-inclusion command or drawing code, so no image is delivered and the package declares no asset dependency.  Editorial formatting only, in addition: an AMS \texttt{amsart} preamble that restates the article's own declarations and macros used by the retained lines, namely the environments \texttt{proposition} on separate counters, together with the standard packages those lines need.  The retained environments print in this document as Proposition 1 in order of appearance, whereas the article prints the same items with the numbering of its own full document.  The complete native source of the article as distributed in the arXiv e-print for this exact version is bundled unchanged as \texttt{sources/191438/source0.tex}.
\begin{proposition} \label{Proposition showing linear independence of alternating sum vectors}
Let $n\ge 3$. The set 
%%%%
$$\mathcal V_n = \{\psi_T\,|\, T\subset V_n, \, 1\le |T|\le n\}$$
%%%%% 
of $2^n-1$ vectors of dimension $2^n-2$, has the property that if we remove any one vector from it, the remaining $2^n-2$ vectors are linearly independent, and thus a bassis for $\mathbb R^{2^n-2}$. In particular, its subset 
%%%%%%%
$$\mathcal T_n = \{\psi_T\,|\, T\subset V_n, 3\le|T|\}$$
%%%%%%%
is linearly independent for all $n\ge 3$. 
\end{proposition}

\begin{proposition} \label{Proposition with the basis for the kernel}
Let as before $\mathcal T_n = \{\psi_T\,|\, T\subset V_n, 3\le|T|\}$ be the set of $2^n-1-{n\choose 2}-n$ elements in the kernel of $S$. Then 
%%%
$$\mathcal B_n = \mathcal T_n \cup \{\varphi_i\}_{i=1,\dots, n-1}$$
%%%%
is a basis for the kernel of $S$.  
\end{proposition}

\begin{proof}
The cardinality of $\mathcal B_n$ agrees with the dimension of $\kerr(S)$, thus we only need to demonstrate the linear independence of $\mathcal B_n$. 

For an index $k\in V_n$ consider the linear functional $P_k:\mathbb R^{2^n-2}\to \mathbb R$ given as 
%%%
$$P_k\left( \sum_{\emptyset\subsetneq C \subsetneq V_n} \lambda _C\, \delta_C\right) = \sum \lambda _C \cdot I(k\in C). $$
%%%
Here $I(k\in C)$ is the indicator function, returning 1 if $k\in C$ and returning $0$ if $k\notin C$. It is easy to see that $P_k$ is linear as it is the composition of an orthogonal projection (onto the subspace of basis vectors $\delta_C$ with $k\in C$) with the trace functional. 

For any $\psi_T\in \mathcal T_n$ and for any choice of $k\in V_n$, we obtain $P_k(\psi_T) = 0$, as we demonstrate next: If $k\notin T$ then clearly $P_k(\psi_T)=0$. So assume that $k\in T$, then 
%%%%%%
\begin{align*}
P_k(\psi_T) & = P_k\left(\sum _{C\subset T} (-1)^{|C|} \delta_C  \right), \cr
& = P_k\left(\sum _{C\subset T-\{k\}} (-1)^{|C|+1} \delta_{C\cup\{k\}}  \right), \cr
& = \sum _{C\subset T-\{k\}} (-1)^{|C|+1},   \cr
& = -\sum _{i=0}^{|T|-1} (-1)^i, \cr 
& = - (1-1)^{|T|-1},\cr
& = 0. 
\end{align*}  
%%%%%%
On the other hand, 
%%%%
$$P_k(\varphi_j) = \left\{ 
\begin{array}{rl}
1 & \quad ; \quad j=k, \cr
-1 & \quad ; \quad j\ne k.
\end{array} 
\right. 
$$
%%%%%%
From this we define new linear functionals $L_k:\mathbb R^{2^n-2}\to \mathbb R$ by setting $L_k = P_k - P_{k+1}$, for $k=1,\dots, n-1$. Then it is still true that $L_k(\psi_T)=0$ for every $\psi _T\in \mathcal T_n$. On the other hand,   
%%%%
\begin{align*}
L_{n-1}(\varphi_j) & = \left\{ 
\begin{array}{rl}
2 & \quad ; \quad j=n-1, \cr
0 & \quad ; \quad j< n-1,.
\end{array} 
\right. \cr
& \cr 
L_{k}(\varphi_j) & = \left\{ 
\begin{array}{rl}
2 & \quad ; \quad j=k, \cr
-2 & \quad ; \quad j=k+1, \cr
0 & \quad ; \quad j\ne k, k+1,
\end{array} 
\right.
\end{align*}
%%%%%%
for $k<n-1$. With this understood, suppose that the set $\mathcal B_n$ were linearly dependent and consider a linear dependence relation:
%%%%%
$$\sum _{i=1}^{n-1}\mu_i\, \varphi_i + \sum _{\psi_T\in \mathcal T_n}\lambda _T\,  \psi_T = 0. $$
%%%%%
Apply firstly the linear functional $P_{n-1}$ to both sides to obtain $\mu_{n-1} = 0$. Now consecutively applying $P_{n-2}, \dots, P_1$ to the linear dependence relation further yields $\mu_{n-2} =0, \dots, \mu_1=0$. Since we already showed that $\mathcal T_n$ is linearly independent (cf. Proposition \ref{Proposition showing linear independence of alternating sum vectors}), it follows that $\lambda _T=0$ for all $T$, proving the proposition. 
\end{proof}

\end{document}
